\section{Disk-Backed Growable Vectors with Stable Virtual Addresses}
\label{sec:disk_backed_growable_vectors}

High-resolution Mandelbrot rendering and reference-orbit techniques (\cref{sec:ref-orbit-calc}) impose
extreme demands on memory management. At deep zoom levels, a single authoritative
reference orbit may contain tens of billions of elements; for example, an orbit
with $2.7\times10^{10}$ iterations already implies a logical memory footprint far
beyond what can be held resident in RAM, even with compact per-iteration storage.
At this scale, both the cost of allocation and the cost of relocation become
dominant design constraints.

Several subsystems in \FractalShark{}—including orbit logging, metadata streams,
debug outputs, and the custom heap allocator—require \emph{growable} storage
while simultaneously maintaining \emph{stable virtual addresses} for previously
returned pointers. This requirement arises naturally when the program hands out
interior pointers into a buffer (e.g., into an orbit record array or a heap
arena). Any reallocation-and-copy strategy would invalidate those pointers and
would require copying enormous amounts of data, which is computationally
infeasible at multi-billion-element scales.

In addition to scale and pointer stability, persistence is a first-class
concern. Reference orbits are expensive to compute but comparatively cheap to
reuse. If orbit data can be left resident on disk and later reloaded simply by
mapping the same backing file back into the process address space, recomputation
can be avoided entirely. This enables efficient multi-pass rendering workflows,
interactive debugging, and reuse of authoritative reference orbits across program
invocations.

To address these requirements, we implement a
\code{GrowableVector<EltT>} abstraction that combines (i) a fixed,
virtually-contiguous address range reserved up front and (ii) incremental growth
within that range by committing additional pages on demand. Growth is performed
in place, so the base address never changes and no copying is required as
capacity increases. The abstraction supports two complementary backing modes:

\begin{enumerate}
  \item \textbf{Anonymous (pagefile-backed) growth.} A large virtual region is
  reserved once via \code{VirtualAlloc(..., MEM\_RESERVE, ...)} and pages are
  subsequently committed within that region using
  \code{VirtualAlloc(base, bytes, MEM\_COMMIT, ...)}. Because each commit
  targets the original base address, successful growth preserves
  \code{m\_Data} and does not require relocation.

  \item \textbf{File-backed (disk-backed) growth.} The reserved address range is
  backed by a disk file and mapped into the process using a section object. In
  this mode, pages are demand-paged from the backing file and become resident
  only as accessed, substantially reducing peak commit requirements when the
  logical capacity greatly exceeds physical memory. When persistence is enabled,
  the same backing file can later be remapped to recover the stored data without
  recomputation.
\end{enumerate}

\subsection{Stable-address growth via reserve-and-commit}
\label{sec:reserve_commit}

\begin{figure}[!htbp]
\centering
\begin{tikzpicture}[
    >=Stealth,
    cell/.style={draw, minimum width=0.55cm, minimum height=0.7cm, font=\tiny},
    committed/.style={cell, fill=blue!20},
    reserved/.style={cell, fill=gray!10, pattern=north east lines, pattern color=gray!30},
    label/.style={font=\scriptsize},
]
  % Reserved region
  \foreach \i in {0,...,19} {
    \ifnum\i<12
      \node[committed] (c\i) at (\i*0.6, 0) {};
    \else
      \node[reserved] (c\i) at (\i*0.6, 0) {};
    \fi
  }

  % Base pointer
  \draw[<-, thick] (c0.south) -- ++(0, -0.6) node[below, label] {\code{m\_Data} (stable)};

  % Capacity marker
  \draw[<->, thick, blue!70!black] (c0.north west) ++(0,0.15) -- ++(7.2, 0)
    node[midway, above, label, blue!70!black] {Committed pages (in use)};

  % Reserved marker
  \draw[<->, thick, gray] (c12.north west) ++(0,0.15) -- ++(4.8, 0)
    node[midway, above, label, gray] {Reserved (uncommitted)};

  % Growth arrow
  \draw[->, thick, orange, dashed] (c11.east) ++(0.1, -0.35) -- ++(1.4, 0)
    node[right, label, orange] {Growth: commit more pages};

  % Total region bracket
  \draw[decorate, decoration={brace, amplitude=5pt, mirror}, thick]
    (c0.south west) ++(0, -1.2) -- ++(12.0, 0)
    node[midway, below=5pt, label] {Virtual address reservation (no reallocation)};
\end{tikzpicture}
\caption[\code{GrowableVector} reserve-and-commit memory layout]{\code{GrowableVector} reserve-and-commit memory layout.
  A large virtual region is reserved at initialization, establishing a
  stable base pointer \code{m\_Data}.  Pages are committed incrementally
  as the vector grows (blue, left).  The reserved-but-uncommitted
  region (hatched, right) consumes no physical memory until needed.
  Because the base address never changes, all previously returned
  pointers remain valid across growth operations.}
\label{fig:growable-vector}
\end{figure}

At initialization time, the vector establishes an address anchor:
\begin{itemize}
  \item In anonymous mode, the anchor is created by reserving a large region
  (typically sized to a fraction of physical memory, or an explicit override)
  and recording the returned base pointer \code{m\_Data}.
  \item In file-backed mode, the anchor is created by mapping a section at a
  base address and thereafter requiring all remaps/extends to preserve that base.
\end{itemize}

Growth proceeds by increasing the logical capacity \code{m\_CapacityInElts}
without moving data. In anonymous mode, this corresponds to committing
additional pages within the reserved region. In file-backed mode, it corresponds
to extending the underlying section and allowing the view to cover the expanded
range. Critically, in both cases we \emph{never} allocate a second buffer and
copy: the address stability is enforced by explicitly requesting the existing
base address and treating any base-address change as a hard error.

\subsection{File-backed mapping for low-commit large buffers}
\label{sec:file_backed_mapping}

When a filename is provided (or a temporary filename is generated), the vector
opens or creates a backing file and constructs a section object over it. We use
native NT section APIs (e.g., \code{NtCreateSection},
\code{NtMapViewOfSection}, and \code{NtExtendSection}) to obtain a
reserve-style mapping suitable for later extension.

Conceptually, the pipeline is:
\begin{enumerate}
  \item \textbf{Open/Create backing file.} The file can be persistent (saved
  results) or temporary (e.g., created with delete-on-close semantics when the
  data is not meant to be retained).
  \item \textbf{Create section.} A section is created with read/write
  protection. For growable vectors, the section is configured so that the view
  can be reserved and later extended.
  \item \textbf{Map a view to establish the base address.} The first mapping
  produces the stable anchor pointer \code{m\_Data}. Subsequent growth must
  preserve this base.
  \item \textbf{Extend on demand.} When capacity increases, the section is
  extended to the new byte size. Because the base mapping remains fixed, all
  previously returned pointers remain valid.
\end{enumerate}

This approach has two key practical benefits:
\begin{itemize}
  \item \textbf{Stable pointers with large logical capacity.} The program may
  reserve and grow to very large capacities (hundreds of GiB or more) without
  the pointer invalidation inherent in \code{realloc}-style strategies.
  \item \textbf{Reduced peak commit pressure.} The OS can keep most of the
  logical buffer non-resident until touched; pages are faulted in as accessed
  and can be evicted back to disk under memory pressure. This is particularly
  effective when the workload has a large addressable dataset but a much smaller
  active working set.
\end{itemize}

\subsection{Trimming and persistence policies}
\label{sec:trim_persist}

The vector supports multiple policies that control whether data is persisted:
\begin{itemize}
  \item \textbf{Persistent save (\code{EnableWithSave}).} Upon \code{Trim()},
  the backing file is truncated to match the used size, ensuring the on-disk
  representation reflects the logical content exactly.
  \item \textbf{Temporary file-backed (\code{EnableWithoutSave}).} Upon
  \code{Trim()}, the view can be remapped to match the used size while
  \emph{reusing the same base address}. This reduces the mapped span without
  sacrificing pointer stability, and the file is automatically deleted when the
  handle closes.
  \item \textbf{Anonymous (\code{DontSave}).} Trimming is a logical operation
  (used size changes) while the reserved region remains as the stable anchor.
\end{itemize}

\subsection{Using the same mechanism for the heap arena}
\label{sec:heap_uses_growable_vector}

The custom heap allocator builds its arena on top of the same
\code{GrowableVector<uint8\_t>} abstraction. The heap is initialized by
creating a file-backed growable vector for the heap arena and committing an
initial chunk of memory. When the allocator needs more space, it expands the
arena by growing the underlying vector in-place and then updating the heap's
wilderness block (coalescing if possible). Because the arena grows within a
reserved and stable virtual range, all heap pointers returned to the program
remain valid across heap expansions, and expansion does not require copying or
relocating the heap.

\subsection{Portability and backing-file considerations}
\label{sec:disk_backed_portability_notes}

The implementation described above relies on Windows-specific virtual memory and
section-mapping semantics, in particular the ability to reserve a virtually
contiguous address range and to grow into that range while preserving a fixed
base address. On Windows, this is achieved using native section objects backed by
ordinary (non-sparse) files whose size is explicitly extended as the logical
capacity grows.

An important consequence of this design choice is that the backing file itself
physically grows on disk as the vector expands, even if large portions of the
mapped address space are never touched. Windows mitigates the runtime memory
impact by demand-paging pages into RAM and evicting them back to disk under
memory pressure, but the on-disk footprint still reflects the full logical extent
of the section. This behavior is acceptable—and often desirable—on Windows,
where file-backed sections are tightly integrated with the operating system's
paging and commit model.

More specifically, Windows treats a section backed by a regular (non-sparse)
file with a well-defined size as a complete and deterministic backing store for
any page in the mapped range. Page faults are resolved by reading from the file,
and evicted pages can be written back to the same file without requiring
additional allocation or metadata decisions at fault time. By explicitly
growing the backing file as capacity increases, the implementation ensures that
the backing storage for every potential page already exists. This aligns with
Windows-native paging expectations and simplifies correctness guarantees around
section growth, page-in, and page-out behavior.

On other platforms, or in alternative designs, a \emph{large sparse file} may be
a more appropriate backing store for growable, disk-backed vectors. Sparse files
allow the logical file size to greatly exceed the actual disk space consumed,
allocating physical blocks only for regions that are written. When combined with
memory-mapped I/O, this can preserve the core advantages of the approach—stable
virtual addresses and low resident memory usage—while also minimizing disk space
consumption for unused regions.

The present implementation intentionally does not rely on sparse-file behavior,
favoring predictability and robustness of paging semantics over minimizing
nominal disk usage. This choice, while well aligned with Windows-native behavior,
has implications for portability: environments that do not support extending a
mapped view in place, or that impose stricter constraints on file-backed mappings
(such as Wine or certain Unix-like kernels), may require alternative strategies.
Possible approaches include explicit sparse-file usage, preallocation of large
sparse sections, or fallback allocators that relax the stable-address requirement
in exchange for broader compatibility.


